\documentclass[12pt]{jlreq}
\usepackage{amsmath, amssymb}

\newcommand{\C}{\mathbb{C}}
\newcommand{\R}{\mathbb{R}}
\newcommand{\Z}{\mathbb{Z}}
\newcommand{\N}{\mathbb{N}}
\newcommand{\F}{\mathbb{F}}
\newcommand{\Prob}{\mathbb{P}}
\newcommand{\Exp}{\mathbb{E}}
\newcommand{\dd}{\,d}
\newcommand{\Ker}{\operatorname{Ker}}
\newcommand{\Impart}{\operatorname{Im}}
\newcommand{\Hom}{\operatorname{Hom}}

\begin{document}

次で定義される遷移系 \((Q,\to)\) を考える。状態集合 \(Q\) は，\(x,y,z\in\{-1,0,1\}\) であり，\(3\) 成分のうちちょうど一つが \(0\) であるような三つ組 \((x,y,z)\) 全体の集合である。遷移に関しては，\(x,y,z\in\{\pm1\}\) として，\(xyz=1\) のときと \(xyz=-1\) のときに応じて，それぞれ，
\[
  (0,y,z)\to(x,y,0)\to(x,0,z)\to(0,y,z)\qquad (xyz=1),
\]
\[
  (x,y,0)\to(0,y,z)\to(x,0,z)\to(x,y,0)\qquad (xyz=-1)
\]
であるとする。これらで定まる \(24\) 通り以外には，遷移は存在しないとする。無限遷移列
\[
  \pi(0)\to\pi(1)\to\pi(2)\to\cdots
\]
のうちで，最初の状態 \(\pi(0)\) が \(q\in Q\) に等しいもの全体を \(\Pi(q)\) と書くことにする。集合 \(X\subseteq Q\) に対して，
\[
  EG(X)=\{q\in Q\mid \exists\pi\in\Pi(q)\ \forall n\ge0\ \pi(n)\in X\},
\]
\[
  AF(X)=\{q\in Q\mid \forall\pi\in\Pi(q)\ \exists n\ge0\ \pi(n)\in X\}
\]
と定める。また \(|X|\) で \(X\) の元の個数を表すとする。

\begin{enumerate}
  \item \(\max\{|X|\mid EG(X)=\emptyset\}\) を求めよ。
  \item \(\max\{|X|\mid EG(AF(X))=\emptyset\}\) を求めよ。
\end{enumerate}

ここで遷移系 \((Q,\to)\) とは有限集合の組であって，\(\to\subseteq Q\times Q\) をみたすものとする。また，\((q,q')\in\to\) のかわりに \(q\to q'\) と書く。

\end{document}